% ----------------------
\subsection{The universality of the gospel}
% ----------------------

	In the introductory essay to this series, I mentioned that future essays would probably range very widely over many different topics. The justification for this, I said, is that the gospel is universal, and so we should expect it to have very broad reach---it influences, and is influenced by, many things. In this second post, I want to talk a little more about that universality, in order to preview some future topics I intend to cover.

	The universality of the gospel of Jesus Christ is neither appreciated nor even recognized by most members of the LDS church. Of course, if you asked them, those members would probably say that the gospel applies everywhere in the universe. But affirming that is one thing, while seeing and accepting its extremely significant consequences is something else entirely.

	Before getting to some of those consequences, let's state more clearly what we mean by the ``universality of the gospel.'' Here is a conditional statement that I consider true:

	\begin{quote}
		\textit{If the gospel of Jesus Christ is true, then it is true throughout the entire universe.}
	\end{quote}
		
	We can ignore the temporal aspects of ``the entire universe'' for now; I'll probably have to take that up in a future essay. All I mean in the consequent, or the ``then'' half of the conditional statement, is that gospel laws apply everywhere in the universe, including to intelligent beings wherever they exist, either on Earth or elsewhere. Again, I think that many members of the LDS church would accept that idea as true.

	(You can already see the extensive ground we need to cover in thinking through the implications of the gospel---we're going to need posts about gospel laws, intelligent beings, other planets, and many other things besides, since even saying ``The gospel being true \textit{means} its laws apply everywhere, to all intelligent beings'' is a philosophical claim which we need to explain and justify.)

	Accepting the conditional statement above as true has an interesting implication. One of the basic rules of logic for conditional statements is that, if the conditional ``If P, then Q'' is true, then the conditional ``If not-Q, then not-P'' is true. Transforming the first conditional into this second one proceeds by a rule called ``transposition.'' So if we apply the rule of transposition to the conditional I gave above, we get a second conditional:

	\begin{quote}
		\textit{If the gospel of Jesus Christ is not true throughout the entire universe, then it is not true.}
	\end{quote}
		
	You can see that all we've done is to switch the parts of the first conditional, then negate both parts. The resulting conditional might be harder for some people to accept, however, even though it is logically equivalent to the first. This conditional says that if the gospel isn't true universally, then it isn't true at all. Another way of putting the point is to say that the gospel being true throughout the universe is a necessary condition on the gospel being true at all.

	I say that this new conditional might be harder to accept than the first because it more clearly reveals what the first commits you to. You're committing yourself to the position that, if there is a place in the universe where the gospel isn't true, then the gospel must be false. But I think accepting these conditionals creates a stronger and more interesting philosophical position than you would get from any of the alternatives. One such alternative might be this position: ``While the gospel may be true here on Earth, it could be false elsewhere.'' Although this view gives you more flexibility than true universality does, I think it's too weak to be desirable or to do justice to claims in the scriptures and from prophets.

	So there are two things to note so far about the universal nature of the gospel. The first is that, as I understand ``universal'' here, it means that the gospel is either true everywhere, or not true at all. Second, on a simple understanding, to say that the gospel is universal means that gospel laws apply everywhere, to all intelligent beings, wherever they are found. (In a later essay, I will introduce the notion of an ``implementation'' of the gospel on a world, which will develop and sharpen this simple understanding of gospel universality.)

	Another point related to the universality of the gospel is the manner in which intelligent beings emerge, or come to be in the first place. The essays in this series are going to take for granted that the theory of organic evolution by natural selection is true. This means that human bodies and, more importantly, human nervous systems, are the products of evolutionary processes. It means further that the bodies and psychological capacities of at least some other intelligent beings, wherever they exist, are also the products of natural evolutionary processes.%
		\footnote{%
			\label{CTR-0002-natural-processes}%
			I've phrased it with ``at least some other'' in order to leave open the possibility that some intelligent beings come from design processes overseen by other intelligent beings. Conscious robots would be an example.
			}

	I will have a great deal to say about evolution later on. I bring it up here only in connection with the gospel's universality. No scientific theory or discovery should have a greater impact on one's thinking about the restored gospel than evolution by natural selection. Many members of the LDS church, and people of faith in general, find it difficult to reconcile this theory with religious belief. But we should accept it without reservation because, perhaps paradoxically, it opens the way to abundant explanatory resources that we do not have access to otherwise.

	More to the point here, evolution is the process by which bodies and, for some organisms, nervous systems come to be on Earth and on other planets. Assuming the universality of the gospel, this fact has extremely significant consequences for how we should understand revelation, righteousness, progression, and God's work in general.

	Exploring those consequences will be an important part of these essays. To cover all that ground, we're going to have to range through cosmology, physics, biology, psychology, and other fields of science, in addition to reflecting on scientific knowledge in general. We'll also be looking at the scriptures and the words of prophets and apostles. To name just a few of the topics, we'll be doing detailed analyses of gospel laws, the universal extent of the gospel, gospel implementations, eternal gender, and many other issues. The outcomes of these analyses, along with their relations to our scientific understanding of the world, will be an opportunity to expand our minds and rejoice in the great plan of salvation that God has revealed.

	In the temple we learn that all truth can be circumscribed into one great whole. In that spirit, a full picture of the gospel---at least, as full a picture as we can achieve---only comes through careful examination of many different truths in many different domains. There is no other way to climb the rainbows.